\chapter{Untergruppen, Erzeugnisse und Nebenklassen}

\section{Untergruppen und Erzeugnisse}

\begin{definition}[Untergruppe und erzeugte Untergruppe]
Eine Teilmenge \(H\subseteq G\) heißt \emph{Untergruppe}, geschrieben
\(H\leq G\), wenn \(1\in H\), \(ab\in H\) und \(a^{-1}\in H\)
für alle \(a,b\in H\) gelten.
Für \(A\subseteq G\) ist die von \(A\) \emph{erzeugte Untergruppe}
\[
 \langle A\rangle=\bigcap\{H\leq G:A\subseteq H\}.
\]
Man nennt \(A\) ein \emph{Erzeugendensystem} von \(G\), wenn
\(\langle A\rangle=G\), und \(G\) \emph{zyklisch}, wenn ein einziges
Element genügt. Es gilt \(\langle\varnothing\rangle=\{1\}\).
\end{definition}

\begin{theorem}[Untergruppenkriterium]
Für \(H\subseteq G\) sind äquivalent:
\begin{enumerate}
\item \(H\leq G\).
\item \(H\ne\varnothing\), und \(ab^{-1}\in H\) für alle \(a,b\in H\).
\item \(H\) ist mit der eingeschränkten Verknüpfung selbst eine Gruppe.
\end{enumerate}
Im dritten Fall stimmen die Einheit und die Inversen von \(H\)
mit denjenigen von \(G\) überein. \cite{milne}
\end{theorem}

\begin{theorem}[Durchschnitte und aufsteigende Vereinigungen]
Der Durchschnitt jeder Familie von Untergruppen von \(G\) ist eine
Untergruppe; der leere Durchschnitt wird als \(G\) verstanden.
Ist \((H_i)_{i\in I}\) eine nichtleere Familie von Untergruppen und
gibt es zu je zwei Indizes \(i,j\) ein \(k\) mit
\(H_i\cup H_j\subseteq H_k\), so ist auch \(\bigcup_{i\in I}H_i\leq G\).
\cite{milne}
\end{theorem}

\begin{theorem}[Wortbeschreibung des Erzeugnisses]
Für jedes \(A\subseteq G\) gilt
\[
 \langle A\rangle
 =\{a_1^{\varepsilon_1}\cdots a_n^{\varepsilon_n}:
 n\geq0,\ a_i\in A,\ \varepsilon_i\in\{1,-1\}\}.
\]
Das Produkt für \(n=0\) bedeutet \(1\).
Es ist die einzige Untergruppe, die \(A\) enthält und in jeder
Untergruppe mit dieser Eigenschaft enthalten ist. \cite{milne}
\end{theorem}

\begin{theorem}[Zyklische Untergruppe eines beliebigen Elements]
Für \(a\in G\) gilt \(\langle a\rangle=\{a^n:n\in\mathbb Z\}\).
Bei unendlicher Elementordnung ist \(n\mapsto a^n\) ein Isomorphismus
\((\mathbb Z,+)\to\langle a\rangle\).
Bei \(\operatorname{ord}(a)=d<\infty\) ist
\[
 \mathbb Z/d\mathbb Z\longrightarrow\langle a\rangle,\qquad
 n+d\mathbb Z\longmapsto a^n
\]
ein Isomorphismus. Diese Alternative gilt für jedes Element jeder Gruppe.
\cite{milne}
\end{theorem}

\section{Produkte von Untergruppen}

\begin{definition}[Mengenprodukt und inverse Menge]
Für \(A,B\subseteq G\) seien
\[
 AB=\{ab:a\in A,\ b\in B\},\qquad A^{-1}=\{a^{-1}:a\in A\}.
\]
Für \(g\in G\) schreibt man \(gA=\{g\}A\), \(Ag=A\{g\}\)
und \(gAg^{-1}=\{gag^{-1}:a\in A\}\).
\end{definition}

\begin{theorem}[Wann ein Untergruppenprodukt eine Untergruppe ist]
Für \(H,K\leq G\) gilt
\[
 HK\leq G\iff HK=KH.
\]
In diesem Fall ist \(HK=\langle H\cup K\rangle\).
Ohne die Bedingung \(HK=KH\) ist das Erzeugnis die Menge aller
endlichen Produkte mit Faktoren aus \(H\cup K\). \cite{milne}
\end{theorem}

\begin{theorem}[Vereinigung zweier Untergruppen]
Für \(H,K\leq G\) ist \(H\cup K\) genau dann eine Untergruppe, wenn
\(H\subseteq K\) oder \(K\subseteq H\). \cite{milne}
\end{theorem}

\section{Nebenklassen und Index}

\begin{definition}[Nebenklassen und Index]
Für \(H\leq G\) heißen \(gH\) die \emph{Linksnebenklasse} und
\(Hg\) die \emph{Rechtsnebenklasse} von \(g\).
Die Menge der Linksnebenklassen heißt \(G/H\); diese Schreibweise
bezeichnet zunächst nur eine Menge, keine Quotientengruppe.
Der \emph{Index} ist \([G:H]=|G/H|\) und darf unendlich sein.
\end{definition}

\begin{theorem}[Nebenklassenkriterien]
Für \(g,k\in G\) und \(H\leq G\) gelten
\[
 \begin{aligned}
 gH=kH&\iff g^{-1}k\in H
       \iff gH\cap kH\ne\varnothing,\\
 Hg=Hk&\iff gk^{-1}\in H
       \iff Hg\cap Hk\ne\varnothing.
 \end{aligned}
\]
Die Linksnebenklassen zerlegen \(G\); dasselbe gilt rechts.
Jede Nebenklasse hat die Mächtigkeit \(|H|\).
Die Zuordnung \(gH\mapsto Hg^{-1}\) verknüpft linke und rechte
Nebenklassen bijektiv. \cite{milne}
\end{theorem}

\begin{theorem}[Lagrange in Kardinalform und Indexformel]
Für jede Untergruppe \(H\leq G\) gilt die Kardinalgleichung
\[
 |G|=[G:H]\cdot|H|.
\]
Für \(K\leq H\leq G\) gilt, auch bei unendlichen Indizes,
\[
 [G:K]=[G:H]\cdot[H:K].
\]
Die Produkte sind Kardinalprodukte; insbesondere werden hier keine
endlichen Gruppen vorausgesetzt. \cite{milne}
\end{theorem}

\begin{theorem}[Index eines Durchschnitts]
Für \(H,K\leq G\) gilt
\[
 [H:H\cap K]\leq[G:K],\qquad
 [G:H\cap K]\leq[G:H]\,[G:K].
\]
Ist \(HK\) eine Untergruppe, so gilt die genauere Gleichung
\[
 [H:H\cap K]=[HK:K].
\]
Alle Indizes werden als Kardinalzahlen verstanden. \cite{milne}
\end{theorem}

\begin{definition}[Doppelnebenklasse]
Für \(H,K\leq G\) und \(g\in G\) heißt
\(HgK=\{hgk:h\in H,\ k\in K\}\) eine \emph{Doppelnebenklasse}.
Die Menge dieser Teilmengen wird mit \(H\backslash G/K\) bezeichnet.
\end{definition}

\begin{theorem}[Doppelnebenklassenzerlegung]
Die Doppelnebenklassen \(HgK\) zerlegen \(G\).
Die in \(HgK\) enthaltenen Linksnebenklassen von \(K\) sind genau die
Mengen \(hgK\) mit \(h\in H\); ihre Anzahl ist
\[
 [H:H\cap gKg^{-1}].
\]
Insbesondere gilt
\[
 |HgK|=[H:H\cap gKg^{-1}]\cdot|K|.
\]
\cite{milne,conradconj}
\end{theorem}
